%\makeglossaries
\makenoidxglossaries

% Danh mục thuật ngữ và từ viết tắt
\newglossaryentry{api}{
    type=\acronymtype,
    name={API},
    description={Giao diện lập trình ứng dụng},
    first={Giao diện lập trình ứng dụng (Application Programming Interface -- API)}
}

\newglossaryentry{rest}{
    type=\acronymtype,
    name={REST},
    description={Kiểu kiến trúc truyền trạng thái đại diện},
    first={Kiểu kiến trúc truyền trạng thái đại diện (Representational State Transfer -- REST)}
}

\newglossaryentry{sdk}{
    type=\acronymtype,
    name={SDK},
    description={Bộ công cụ phát triển phần mềm},
    first={Bộ công cụ phát triển phần mềm (Software Development Kit -- SDK)}
}

\newglossaryentry{rbac}{
    type=\acronymtype,
    name={RBAC},
    description={Kiểm soát truy cập dựa trên vai trò},
    first={Kiểm soát truy cập dựa trên vai trò (Role-Based Access Control -- RBAC)}
}

\newglossaryentry{jwt}{
    type=\acronymtype,
    name={JWT},
    description={Chuẩn token dùng để truyền thông tin xác thực dạng JSON},
    first={JSON Web Token (JWT)}
}

\newglossaryentry{websocket}{
    type=\acronymtype,
    name={WebSocket},
    description={Giao thức giao tiếp hai chiều giữa máy khách và máy chủ},
    first={WebSocket}
}

\newglossaryentry{webrtc}{
    type=\acronymtype,
    name={WebRTC},
    description={Công nghệ giao tiếp âm thanh, hình ảnh và dữ liệu thời gian thực trên trình duyệt},
    first={Web Real-Time Communication (WebRTC)}
}

\newglossaryentry{html}{
    type=\acronymtype,
    name={HTML},
    description={Ngôn ngữ đánh dấu siêu văn bản},
    first={Ngôn ngữ đánh dấu siêu văn bản (HyperText Markup Language -- HTML)}
}

\newglossaryentry{css}{
    type=\acronymtype,
    name={CSS},
    description={Ngôn ngữ định kiểu giao diện web},
    first={Cascading Style Sheets (CSS)}
}

\newglossaryentry{sql}{
    type=\acronymtype,
    name={SQL},
    description={Ngôn ngữ truy vấn có cấu trúc},
    first={Structured Query Language (SQL)}
}

\newglossaryentry{orm}{
    type=\acronymtype,
    name={ORM},
    description={Kỹ thuật ánh xạ đối tượng trong chương trình với bảng trong cơ sở dữ liệu quan hệ},
    first={Object-Relational Mapping (ORM)}
}

\newglossaryentry{crud}{
    type=\acronymtype,
    name={CRUD},
    description={Nhóm thao tác tạo, đọc, cập nhật và xóa dữ liệu},
    first={Create, Read, Update, Delete (CRUD)}
}

\newglossaryentry{erd}{
    type=\acronymtype,
    name={ERD},
    description={Biểu đồ thực thể -- liên kết},
    first={Entity Relationship Diagram (ERD)}
}

\newglossaryentry{uml}{
    type=\acronymtype,
    name={UML},
    description={Ngôn ngữ mô hình hóa thống nhất},
    first={Unified Modeling Language (UML)}
}

\newglossaryentry{ui}{
    type=\acronymtype,
    name={UI},
    description={Giao diện người dùng},
    first={User Interface (UI)}
}

\newglossaryentry{ux}{
    type=\acronymtype,
    name={UX},
    description={Trải nghiệm người dùng},
    first={User Experience (UX)}
}
